\documentclass[11pt]{article}
\usepackage[margin=1in]{geometry}
\usepackage{amsmath,amssymb,mathtools,booktabs,array}
\usepackage[T1]{fontenc}
\usepackage{lmodern}
\usepackage[hidelinks]{hyperref}
\usepackage{microtype}
\newcommand{\adj}{^{\dagger}}
\newcommand{\norm}[1]{\left\lVert#1\right\rVert_2}
\newcommand{\file}[1]{\nolinkurl{#1}}
\title{MetalClover GPU implementation notes\\\large Schur multigrid, local SVD, and Apple Metal}
\author{Notes from the MetalClover source and recorded experiments}
\date{15 November 2026}
\begin{document}
\maketitle

\paragraph{Scope.}
These notes describe the implementations in this repository, with emphasis on
the sparse three-level multigrid (MG) path and the optional near-null/SVD setup.
The older dense two-level prototype is identified separately. Mathematical
equalities below describe the unrounded operators; device arithmetic and stored
hierarchy coefficients are predominantly FP32. Source names in Section~\ref{sec:source}
provide a guide to the implementation and its recorded validation.
Unless stated otherwise, source and study filenames refer to files under
\file{metal-bsr-prototype/} relative to the repository root. Public API usage
and build instructions are in \file{docs/API.md} and \file{docs/BUILD.md}.

\section{Fine equation and even/odd elimination}
At each lattice site a fermion has three colors and four spins, hence twelve
complex components. Write the clover--Wilson equation as
\begin{equation}
 Mx=b,\qquad M=A+H,
\end{equation}
where $A$ is site-local and $H$ contains nearest-neighbor hopping. In the
isotropic synthetic prototype, $A=(m_0+4)I+C$ and $H=-D\!\!\!\!/\,/2$.
The real-field adapter uses the imported anisotropic action: direction weights,
spatial/temporal clover coefficients, stout gauge state, and temporal boundary
sign are part of that operator. The isotropic formula must not replace those
action-dependent coefficients.

In the implemented gamma basis, $\gamma_5=\operatorname{diag}(I_6,-I_6)$
up to the interleaved color/spin storage ordering, and each $A(x)$ consists
of two Hermitian $6\times6$ blocks. With even sites ordered before odd sites,
\begin{equation}
 M=\begin{pmatrix}A_{ee}&H_{eo}\\H_{oe}&A_{oo}\end{pmatrix},\qquad
 S=A_{oo}-H_{oe}A_{ee}^{-1}H_{eo}.
 \label{eq:schur}
\end{equation}
Here the first subscript is the destination parity. The solver retains the
odd system, with $12V/2$ complex unknowns:
\begin{align}
 Sx_o&=b_o-H_{oe}A_{ee}^{-1}b_e,\\
 x_e&=A_{ee}^{-1}(b_e-H_{eo}x_o).
\end{align}
One application of $S$ therefore needs two parity-changing hops, an even
local solve, an odd local multiplication, and a subtraction. $S$ is applied
without forming its full two-hop matrix. No normal equations are used for
the inversion.

\subsection{Mixed-precision defect correction}
The physical solution and acceptance residual use CPU FP64. At outer step $k$,
\begin{align}
 r^{(k)}&=b-M_{64}x^{(k)},\\
 g_o&=r_o^{(k)}-H_{oe}A_{ee}^{-1}r_e^{(k)},\\
 S_{32}\delta_o&\simeq (g_o)_{32},\\
 \delta_e&=A_{ee}^{-1}(r_e^{(k)}-H_{eo}\delta_o),\\
 x^{(k+1)}&=x^{(k)}+\operatorname{FP64}(\delta).
\end{align}
Reduction, inner solution, and reconstruction are GPU FP32 operations on the
Metal paths. Acceptance requires the full original equation to satisfy
$\norm{b-M_{64}x}/\norm b\leq\varepsilon$ for every RHS. Zero RHS use an
absolute criterion. The early synthetic studies used $10^{-10}$; the cited
real-field studies requested $10^{-12}$. An inner Arnoldi estimate or coarse
residual alone cannot accept a physical solution.

\section{Aggregation, prolongation, and restriction}
Partition full lattice coordinates into disjoint rectangular aggregates
$\Omega_a$. Only odd sites inside each aggregate participate in the fine
transfer. For example, an even-sided $4^4$ aggregate has $128$ odd sites.
The implementation also supports rectangular and odd-sided aggregate shapes
where the lattice partition permits them; parity counts then follow the actual
site lists.

Let $n_v$ be the number of retained vectors \emph{per chirality}. There are
$2n_v$ coarse components per aggregate. For $\chi\in\{+,-\}$, define local
orthonormal columns $p_{a\chi j}$ supported on $\Omega_a$ and chirality $\chi$.
If $q$ denotes a fine color/spin component, the prolongator is
\begin{equation}
 P_{(x,q),(a,\chi,j)}=
 \begin{cases}
 p_{a\chi j}(x,q),&x\in\Omega_a,\quad q\in\chi,\\
 0,&\text{otherwise},
 \end{cases}
 \quad j=1,\ldots,n_v.
 \label{eq:P}
\end{equation}
Thus interpolation is aggregate-local, with $n_v$ nonzero weights per fine
component. It is not geometric linear interpolation between neighboring
aggregate centers. Its action and the restriction are
\begin{align}
 (Pe)(x,q)&=\sum_{j=1}^{n_v}p_{a(x)\chi(q)j}(x,q)e_{a(x)\chi(q)j},\\
 (Rr)_{a\chi j}&=\sum_{x\in\Omega_a}\sum_{q\in\chi}
       \overline{p_{a\chi j}(x,q)}\,r(x,q),\qquad R=P\adj.
 \label{eq:R}
\end{align}
No extra volume normalization or full-weighting stencil appears in $R$.
Local orthonormality and disjoint support give
\begin{equation}
 P\adj P=I,\qquad \gamma_{5,o}P=P\Gamma_5,\qquad
 \Gamma_5\big|_a=\operatorname{diag}(I_{n_v},-I_{n_v}).
\end{equation}
$PP\adj$ is the Euclidean orthogonal projector onto the interpolation space;
it is generally not the identity. These identities hold to setup accuracy
before the FP32 cast.

\section{The coarse operator: sparse Galerkin blocks}
The first coarse matrix is the projection of the \emph{odd Schur complement}:
\begin{equation}
 A_1=P_0\adj S P_0,\qquad
 (A_1)_{a\chi i,b\psi j}
 =\sum_{x,q;y,t}\overline{P_{(x,q),(a,\chi,i)}}
 S_{(x,q),(y,t)}P_{(y,t),(b,\psi,j)}.
 \label{eq:galerkin}
\end{equation}
It has dense complex $2n_v\times2n_v$ blocks on a sparse aggregate graph.
The blocks generally mix both coarse chiralities and all retained vectors.
The coarse operator is not obtained by fitting a new SU(3) link or by simply
reusing the original twelve-component Wilson stencil.

\subsection{Local construction and coarse sparsity}
Let $P_a$ denote the columns supported in aggregate $a$ and let $e$ be an
even site. Define the small matrix
\begin{equation}
 T_{ea}=(H_{eo}P_a)(e),\qquad T_{ea}\in\mathbb C^{12\times2n_v}.
\end{equation}
Using $H_{oe}=\gamma_{5,o}H_{eo}\adj\gamma_{5,e}$ and chirality preservation,
the blocks can be assembled as
\begin{equation}
 (A_1)_{ab}=\delta_{ab}\,P_a\adj A_{oo}P_a
 -\Gamma_5\sum_e T_{ea}\adj\gamma_{5,e}A_{ee}^{-1}T_{eb}.
 \label{eq:localassembly}
\end{equation}
The first term includes a sum over odd sites in $a$. In the second term only
even sites adjacent to odd support in both $a$ and $b$ contribute. Accordingly,
the graph is built by joining the aggregates that touch the same even site,
plus the diagonal. It includes the couplings induced by two fine hops; it
should not be assumed to have only the nearest-neighbor geometric coarse
stencil. Contributions are accumulated in CPU FP64, with row ownership
avoiding assembly atomics. This is an exact algebraic Galerkin construction
up to arithmetic rounding, without global dense column probing.

The representation is block compressed sparse row: row offsets, aggregate
column indices, and contiguous dense blocks. Application computes
\begin{equation}
 (A_1u)_a=\sum_{b\in\mathcal N(a)}(A_1)_{ab}u_b.
\end{equation}
Only diagonal blocks are inverted for block-Jacobi smoothing. They are the
diagonal blocks of the projected matrix, including Schur contributions;
they are not just restrictions of $A_{oo}$.

Since $S\adj=\gamma_{5,o}S\gamma_{5,o}$,
\begin{equation}
 A_1\adj=\Gamma_5A_1\Gamma_5.
\end{equation}
This is gamma5 Hermiticity, not ordinary Hermiticity or positive definiteness.
It explains the use of non-Hermitian Krylov solvers rather than ordinary CG.

\subsection{A second coarsening}
The three-level path aggregates coarse sites and constructs another
chirality-preserving $P_1$, followed by
\begin{equation}
 A_2=P_1\adj A_1P_1,\qquad
 (A_2)_{AB}=\sum_{a\in A,b\in B}(P_1)_a\adj(A_1)_{ab}(P_1)_b.
\end{equation}
This is a sparse block contraction. The second transfer's input chirality
is the coarse $+/-$ component grouping, rather than physical spin indices.
The near-null/SVD option changes the first transfer; the documented second
transfer retains the filtered-vector/local-orthogonalization construction.

\section{Building the interpolation space}
\subsection{Earlier filtered-vector setup}
The original setup starts from deterministic constant-plus-noise complex
test vectors and filters them with homogeneous damped clover-block Jacobi:
\begin{equation}
 v\leftarrow (I-\omega_{\rm setup}A_{oo}^{-1}S)v.
\end{equation}
Aggregate/chirality restrictions are then orthonormalized with two-pass local
Gram--Schmidt. The corresponding second-level filter uses the coarse diagonal
block inverse. The setup weight and number of sweeps are separate parameters
from the pre/post-smoother settings during a solve. The original prototype
filters in CPU FP64; the real-field path can filter fine vectors on the GPU
in FP32 and perform local orthogonalization in CPU FP64.

\subsection{Homogeneous-solve near-null candidates}
The optional builder in \file{near_null_svd.hpp} uses a larger candidate pool.
For each deterministic complex Z4 start, independent of physical RHS, it:
\begin{enumerate}
 \item Optionally warms $y$ by damped homogeneous Jacobi sweeps.
 \item Forms $f=Sy$ and approximately solves $Sz=f$ from zero with GPU FP32
 FGMRES and local clover-block preconditioning.
 \item Forms the remaining error $v=y-z$ and explicitly measures $Sv$.
 \item Normalizes the full candidate: $\widehat v=v/\norm v$, using an FP64 norm.
\end{enumerate}
The key identity is
\begin{equation}
 Sv=f-Sz.
\end{equation}
An incomplete solve leaves error that the chosen approximate inverse removes
poorly. Such errors provide useful coarse-space candidates. The scale-invariant
quantity $\norm{Sv}/\norm v$ indicates how near-null a candidate actually is;
a small residual relative to $\norm{Sy}$ alone does not prove that property.
An exact solve of an invertible $S$ would leave $v=0$, so finite inexact setup
is intentional. The implementation permits a setup iteration cap and records
the achieved reductions. All 24 setup candidates in the recorded large-volume
study reached their 128-iteration cap rather than the requested $0.001$ target.

\subsection{What the SVD does}
For each aggregate/chirality gather the $m$ normalized candidates into
\begin{equation}
 V_{a\chi}=[\widehat v_1|_{a\chi}\ \cdots\ \widehat v_m|_{a\chi}]
   =U\Sigma W\adj,\qquad \sigma_1\geq\cdots\geq\sigma_m\geq0.
\end{equation}
The code calls Accelerate LAPACK \file{zgesvd} in CPU FP64, retaining
\begin{equation}
 [p_{a\chi1},\ldots,p_{a\chi n_v}]=U(:,1:n_v).
\end{equation}
These are the leading \emph{left} singular vectors of a local sample matrix.
They are not singular vectors computed directly from the global Dirac matrix.
Large singular values are retained because they represent the dominant
\emph{candidate covariance}; the candidates already encode difficult errors.
The local optimality statement is
\begin{equation}
 \min_{Q\adj Q=I_{n_v}}\|V-QQ\adj V\|_F^2
 =\sum_{j>n_v}\sigma_j^2,\qquad
 \eta=\frac{\sum_{j\leq n_v}\sigma_j^2}{\sum_j\sigma_j^2}.
\end{equation}
Thus the SVD compresses locally redundant errors into a fixed-size basis.
Energy retention $\eta$ measures representation of these samples, not a
guarantee on the global low spectrum or the solver's convergence rate.
The builder checks local orthogonality and rejects retained numerical rank
with $\sigma_{n_v}\leq10^{-12}\sigma_1$.

\subsection{Bounded-memory incremental compression}
The later implementation can compress a full pool to weighted columns
$U_k\Sigma_k$, append new candidates $B$, and recompress
\begin{equation}
 [\,U_k\Sigma_k\ \ B\,].
\end{equation}
The weights matter because
$(U_k\Sigma_k)(U_k\Sigma_k)\adj=U_k\Sigma_k^2U_k\adj$ preserves the
retained covariance. Carrying $U_k$ alone would assign equal weights to
directions of unequal importance. Final interpolation uses unweighted $U_{n_v}$
so that $P\adj P=I$. Intermediate truncations mean this is generally different
from one SVD of all original candidates.

The $32^3\times256$ study used $m=24$, $n_v=8$ per chirality, a nine-column
weighted carry, and four new columns per batch: a thirteen-column pool.
One FP32 odd vector is $384$ MiB, so the pool occupies $4.875$ GiB instead of
$9$ GiB for all 24 vectors. The recorded energy lower bound relative to the
original normalized candidates is $96.555\%$; it is not the retained fraction
of a full 24-column SVD. The earlier $24^3\times256$ study used one final
all-candidate SVD. These are different setup variants.

\section{Multigrid cycle and Krylov solver}
At level $\ell$, let $B_\ell$ be the local block inverse: $A_{oo}^{-1}$ on
the fine odd level and $\operatorname{blockdiag}(A_\ell)^{-1}$ on coarse levels.
Starting with $z=0$, a cycle for residual $r$ performs
\begin{align}
 z&\leftarrow z+\omega_\ell B_\ell(r-A_\ell z)
       &&\text{for each pre-step},\\
 r_c&=P_\ell\adj(r-A_\ell z),\\
 A_{\ell+1}e_c&\simeq r_c,\\
 z&\leftarrow z+P_\ell e_c,\\
 z&\leftarrow z+\omega_\ell B_\ell(r-A_\ell z)
       &&\text{for each post-step},
\end{align}
where $A_0=S$. Fine smoothing uses the clover block, not the full diagonal
of the two-hop Schur matrix. For an exact coarse solve, the coarse error
propagator is $I-P_\ell A_{\ell+1}^{-1}P_\ell\adj A_\ell$; it annihilates
error in $\operatorname{range}(P_\ell)$. For one pre/post step the full error
propagator is
\begin{equation}
 (I-\omega B A_\ell)
 (I-P_\ell A_{\ell+1}^{-1}P_\ell\adj A_\ell)
 (I-\omega B A_\ell).
\end{equation}
The applied FP32/inexact cycle only approximates this expression.

The old two-level prototype formed a small dense $A_1$ by column probing and
stored a CPU-FP64-computed dense inverse, cast to FP32. It limits the coarse
dimension to 512. The scalable three-level path instead uses an inner coarse
FGMRES solve, preconditioned by another MG cycle, and a bottom block-Jacobi
BiCGStab solve. No global dense coarse inverse is used there. Restart sizes,
targets and iteration caps vary between studies; the original three-level
example used fine FGMRES(8), coarse FGMRES(4), coarse target $0.1$, and an
eight-iteration bottom cap. Its bottom target is a left-preconditioned norm;
the unpreconditioned residual is recorded separately.

Fine FGMRES is right-preconditioned. At step $j$ it stores
\begin{equation}
 z_j=\mathcal B_jv_j,\qquad w_j=Sz_j,\qquad
 SZ_m=V_{m+1}\overline H_m,
\end{equation}
and updates $x=x_0+Z_my$ with
$y=\arg\min_y\|\beta e_1-\overline H_my\|_2$.
Keeping $Z_m$ explicitly permits variable/inexact coarse corrections.
Two-pass Gram--Schmidt is used; small Givens/least-squares control is CPU FP64.
At restart or a predicted early stop, the true FP32 Schur residual is refreshed.
Multiple RHS share coefficients and dispatches but have separate Krylov bases,
Hessenberg problems and convergence masks: this is batched independent Krylov,
not block GMRES.

\section{Apple Metal implementation choices}
\paragraph{Compact hopping and spin structure.}
The compact fine operator stores forward gauge links, neighbor tables, parity
maps, and clover data instead of expanded copies of a fine sparse matrix.
Backward links are obtained by conjugate transpose at the neighbor. Wilson
spin projectors use spin permutations and complex phases, followed by small
color contractions; dense $12\times12$ hopping blocks need not be stored.
Anisotropy and the temporal wrap sign are applied according to the action.

\paragraph{RHS-contiguous layout and specialization.}
Fine and coarse fields use a component-major layout
$\texttt{field[component * RHS + rhs]}$. Adjacent threads can process adjacent
RHS for the same component while reusing common operator data. Shaders use a
specialized \file{RHS} constant. Coarse kernels contract BSR blocks using
FP32 complex arithmetic represented by \file{float2}, with explicit fused
multiply-add operations. The supplied kernels are conventional compute
kernels; these notes do not assume tensor/matrix-unit acceleration.

\paragraph{Local clover compression and LDL reuse.}
Packed Hermitian storage needs six real diagonal entries and fifteen complex
lower entries per chiral block: 36 reals instead of 36 complex numbers.
The low-memory path stores $A=LDL\adj$, with unit-diagonal lower $L$, in the
same 36-real footprint. Multiplication evaluates $L[D(L\adj x)]$; inversion
uses forward substitution, diagonal division and backward substitution.
The same odd factor serves Schur multiplication and smoothing; the even
factor serves elimination/reconstruction. There is no separate odd inverse
array on this path. FP64 setup checks positive local pivots for the tested
action; this factorization path is not an unrestricted indefinite solver.
The recorded large-volume local storage falls from $6.75$ to $2.25$ GiB.
The current kernel computes a six-component local sequence per output thread;
the compact storage saves memory but does not remove all duplicated local work.

\paragraph{Resident workspaces and synchronization.}
Fine Krylov vectors, preconditioned vectors, transfer weights, sparse coarse
blocks and cycle workspaces remain in Metal buffers during inner solves.
Shared-storage buffers exploit Apple unified memory, but still require command
completion before CPU consumption. Dependent GPU dispatches use buffer barriers;
host control waits where small reductions or nested solver control require it.
Only small reduced coefficients need return to CPU during Arnoldi, not the full
Krylov basis. Some operator interfaces still require basis-to-work-buffer copies.
The implementation includes ordinary dispatch sizes such as 128 threads per
threadgroup; this is not evidence of a universally optimal hardware tuning.

\paragraph{Bounded memory and FP64 streaming.}
Restart length and resident RHS tile size bound the large fine-vector stacks.
The public array interface reuses a context/hierarchy and processes GPU tiles
(default cap eight); the largest cited standalone run held only one RHS.
The low-memory driver maps the FP64 solution to a file and streams residual
work in eight-time-slice slabs. CPU FP64 clover blocks are regenerated from
the original FP64 gauge, rather than promoted from rounded GPU factors.
Sources can be regenerated deterministically, and scratch pages are released
between slabs. These choices make full-precision acceptance feasible without
keeping full source, residual and clover arrays resident simultaneously.
Metal allocation and CPU RSS overlap in unified memory and must not be added.

\paragraph{Remaining costs.}
MG reduces fine-operator work but adds transfer, coarse, reduction and host
control costs. Fixed bottom dispatch budgets can include masked/padded work.
Full FP64 residual evaluation and file writeback remain in solve timing.
Reducing iteration counts or increasing nominal operator GFLOP/s does not by
itself establish a faster total solve. Hierarchy construction is a separate
cost amortized when the gauge/action/mass-specific hierarchy is reused.

\section{What is taken from Eloy Romero's mgproton?}
The repository explicitly credits Eloy's \file{mgproton} work in Chroma and
SuperbBLAS as the starting point for the sparse color/spin representation and
the investigation of even/odd reduction, chirality-split aggregation, Galerkin
coarse operators, and Krylov methods. More specifically, the optional
\file{getMGProlongator}-inspired structure is:
\begin{enumerate}
 \item Generate errors of approximate homogeneous solves,
       $z\simeq S^{-1}(Sy)$ and $v=y-z$.
 \item Normalize candidates globally, then gather them by aggregate/chirality.
 \item Compress the local candidate matrices by SVD and use leading left
       singular vectors as interpolation weights.
 \item Optionally carry singular-value-weighted directions between candidate
       batches, then form the final orthonormal transfer.
\end{enumerate}
The shared structure is algorithmic. MetalClover supplies its own dedicated
CPU/Metal kernels, sparse hierarchy assembly/application, memory management,
solver orchestration and mixed-precision checking. These notes do not claim
that the hierarchy is identical to Eloy's, or that \file{mgproton} is called
as the GPU solver. Agreement with original Chroma establishes the physical
operator and accepted solutions, not equality of MG setup or convergence.
This provenance account is supported by the repository's explicit attribution
and comments; no separate upstream \file{mgproton.cc} checkout was examined
for a line-by-line comparison.

\section{Recorded example and implementation map}\label{sec:source}
The near-null/SVD $32^3\times256$, $m_0=-0.0860$ study on a 32 GiB M1 Max
used $4^4$ fine aggregates, eight retained vectors per chirality at both
transfers, and fine restart eight. The first coarse grid has
$(8\cdot8\cdot8\cdot64)$ aggregate sites and 16 components per site,
or $524{,}288$ unknowns, illustrating why a global dense coarse inverse is
inappropriate. Mean warm solve times were $474.499$ s/RHS for near-null/SVD MG
and $2189.472$ s/RHS for BiCGStab ($4.61\times$). One hierarchy cost
$964.986$ s; including it over the two measured RHS gives $956.992$ s/RHS.
Every saved MG field passed original Chroma FP64 checks, with maximum relative
residual $1.7355\times10^{-14}$. These are two single-resident-RHS measurements,
not measured 96- or 192-RHS throughput.

\begin{description}
 \item[Fine operator and compact storage:]\hfill\\
 \file{compact_clover.hpp}, \file{compact_clover.metal},
 \file{low_memory_clover.hpp}, \file{packed_clover.metal}.
 \item[Transfers and sparse hierarchy:]\hfill\\
 \file{three_level.hpp}, \file{three_level.mm},
 \file{sparse_coarse.hpp}, \file{sparse_coarse.metal}.
 \item[Cycle tuning and near-null setup:]\hfill\\
 \file{mg_real_tune.hpp}, \file{mg_real_tune.metal},
 \file{near_null_svd.hpp}.
 \item[Krylov and original dense prototype:]\hfill\\
 \file{fgmres.hpp}, \file{bicgstab.hpp},
 \file{multigrid.hpp}, \file{multigrid.metal}.
 \item[Full checks and integration:]\hfill\\
 \file{low_memory_solve.mm}, \file{compact_residual.hpp},\\
 \file{production_bridge.hpp}, \file{metalclover_api.mm}.
 \item[Evidence and historical configurations:]\hfill\\
 \file{MULTIGRID.md}, \file{THREE_LEVEL.md}, \file{LOW_MEMORY.md},
 \file{NEAR_NULL_SVD.md}, \file{NEAR_NULL_SVD32.md},
 \file{BATCH_RHS.md}, and the root \file{README.md} provenance statement.
\end{description}
The hierarchy tests cover local orthogonality, gamma5 Hermiticity, Galerkin
action identities, transfer/cycle agreement between CPU and Metal, and full
solution residuals. Performance numbers above are the repository's recorded
experiments; preparing these notes does not rerun the numerical studies.
\end{document}
